\clearpage
\section{Implementarea sistemului informațional}

După cum a fost stabilit sistemul librărie online se
va numi \textbf{Book Mall} va avea trei tipuri de actori, și posibilitatea de a 
gestiona utilizatori și cărți. Sistemul respectiv va fi disponibil pe pagina de
GitHub pe adresa \url{https://github.com/catalinplesu/bookmall}, respectiv va 
avea o licență de tip open-source.

\subsection{Selectarea tehnologiei}
Limbajul de programare este mai puțin important când vorbim despre aplicații
web deoarece crearea unei astfel de aplicații ar fi extrem de dificilă fără
utilizarea unui framework specializat. Printre cele mai populare frameworkuri
pentru realizarea aplicațiilor web sunt:
\begin{itemize}
    \item Express(JS);
    \item Django(PY);
    \item Rails(RB);
    \item Larvel(PHP);
    \item Spring(Java);
    \item ASP.NET(C\#).
\end{itemize}
Am decis să utilizez ASP.NET datorită cunoștințelor recent obținute în cadrul
cursului Tehnologii Web, astfel proiectul va fi dezvoltat în limbajul C\# iar
în calitate de IDE va fi folosit Visual Studio. Menționz că proiectul va fi
realizat în ASP.NET Framework ceea ce îl va limita daor pentru windows ceea ce
ține de server.

În adție la acest framework va fi utilizat și Enitity Framework care este un 
framework de tip ORM care permite accesarea unei baze de date SQL fără a scrie 
cod SQL ceea ce va face dezvoltarea mai ușoară.

Proiectul dat va fi dezvoltat urmând arhitectura Buisness Logic ceea ce
înseamnă că aplicația va fi divizată în mai multe straturi, proiectul dat va
conține următoarele subproiecte: 
\begin{itemize}
\item BooMall.BuisnessLogic
\item BooMall.Domain
\item BooMall.Helpers
\item BooMall.Web
\end{itemize}

\subsection{Rezultatul obținut}

În figura \ref{fig:impl_project} poate fi observată structura proiectului respectiv 
în conformitate cu buisness logic și tabelul cu utilizatori.
\begin{figure}[H]
    \begin{center}
        \includegraphics[width=15cm]{impl/project}
        \caption{Structura proiectului, Tabelul cu utilizatori}
        \label{fig:impl_project}
    \end{center}
\end{figure}

La momentul scrierii acestui raport sistemul nu a fost implementat complet însă
a fost implementată o mare parte. În anexa \ref{anex:bookmall} pot fi observate
mai multe pagini din aplicația creată care corespund funcționalului care a fost
implementat cu succes:

\begin{itemize}
\item pagina principală, figura \ref{fig:impl_home};
\item registrarea, figura \ref{fig:impl_signup};
\item logarea, figura \ref{fig:impl_login};
\item încărcarea cărților, figura \ref{fig:impl_create};
\item descărcarea cărților;
\item setările utilizatorilor, figura \ref{fig:impl_settings}.
\end{itemize}

Iar prin funcționalul care necesită implemenat sunt:
\begin{itemize}
\item resetarea parolei;
\item editarea cărților;
\item ștergerea cărților;
\item ștergerea utilizatorilor;
\item cumpărarea cărților;
\item cumpărarea cărților;
\item setările utilizatorilor.
\end{itemize}
